% GDZ PPN236018647,native548lower–549;HKUST549–550.
\NumberedItem{69}
\indent Si
\begingroup\makeatletter\let\veqno\@@leqno\makeatother
\begin{equation*}\tag*{(A)}
x^{\mu+\nu}+ax^{\mu+\nu-1}+bx^{\mu+\nu-2}+\ldotp\,\ldotp\,\ldotp+n
\end{equation*}
\noindent per
\begin{equation*}\tag*{(B)}
x^\mu+\alpha x^{\mu-1}+\beta x^{\mu-2}+\ldotp\,\ldotp\,\ldotp+m
\end{equation*}
\par\endgroup

\clearpage\SourceLeaf{549}\ArithmeticHead{520}{L}{TAGEBUCH MIT ERLÄUTERUNGEN.}
\noindent dividatur atque omnes coefficientes in (A) \(a,b,c,\) etc. sint numeri integri, coefficientes vero omnes in (B) rationales, etiam hi omnes erunt integri ultimumque \(n\) ultimus \(m\) metietur.
\par\noindent\hfill[1797] Iul. 23.

\medskip\begingroup\small
\indent Dieser Satz wird im art. 42. der \textit{Disquisitiones arithmeticae}, Werke I, S. 34 bewiesen. \textsc{Gauss} benutzt ihn auch im art. 341. der \textit{Disquisitiones arithmeticae}, Werke I, S. 418, 419 beim Beweise der Irreduzibilität der Kreisteilungsgleichung für den Fall eines Primzahlgrades. Mithin muß \textsc{Gauss} die Aussage der Tagebuchnotiz Nr. 40, die ja auf den Satz des art. 341. hinausläuft, am 9. Oktober 1796 entweder anders als in den \textit{Disquisitiones arithmeticae} oder noch nicht vollständig bewiesen haben. Vergl. auch die Anmerkung zum art. 42. der \textit{Disquisitiones arithmeticae}, Werke I, S. 475, wo das Datum 1797 Jul. 22. angegeben ist. — In der Handschrift stehen in den Exponenten der Ausdrücke (A), (B) die Buchstaben \(m,n\) statt \(\mu,\nu\).
\par\noindent\hfill\textsc{Klein.\quad Loewy.}
\par\endgroup

\NumberedItem{70}
\indent Forsan omnia producta ex
\[
(a+b\rho+c\rho^2+d\rho^3+\ldotp\,\ldotp\,\ldotp)
\]
\noindent designante \(\rho\) omnes radices prim[itivas] aeq[uationis] \(x^n=1\) ad formam
\[
(x-\rho y)(x-\rho^2y)\ldotp\,\ldotp\,\ldotp
\]
\noindent reduci possunt[\textsuperscript{*})]\begingroup\renewcommand{\thefootnote}{[*)}\footnotetext{In der Handschrift steht \textls[120]{potest}.]}\endgroup. Est enim:
\[
\begin{aligned}
(a+b\rho+c\rho^2)\times(a+b\rho^2+c\rho)&={}\\
&(a-b)^2+(a-b)(c-a)+(c-a)^2
\end{aligned}
\]
\[
\begin{aligned}
(a+b\rho+c\rho^2+d\rho^3)\times(a+b\rho^3+c\rho^2+d\rho)&={}\\
&(a-c)^2+(b-d)^2
\end{aligned}
\]
\[
\begin{aligned}
&(a+b\rho+c\rho^2+d\rho^3+e\rho^4+f\rho^5)\times\qquad=(a+b-d-e)^2\\
&\qquad-(a+b-d-e)(a-c-d-f)+(a-c-d-f)^2\\
&\hspace{15em}=(a+b-d-e)^2\\
&\qquad+(a+b-d-e)(b+c-e-f)+(b+c-e-f)^2.
\end{aligned}
\]
\noindent\hfill Vid. Febr. 4.

\indent Falsum est. Hinc enim sequeretur, e binis numeris in forma Pr[oducti] e \((x-\rho y)\) contentis productum in eadem forma esse, quod facile refutatur.
\par\noindent\hfill[1797 Iul.]

\medskip\begingroup\small
\indent Im zweiten und dritten Gliede der letzten Gleichung scheinen einige Vorzeichen verschrieben zu sein, denn es gilt die Identität:
\[
\begin{aligned}
&-(a+b-d-e)(a-c-d+f)+(a-c-d+f)^2\\
&\qquad=-(a+b-d-e)(b+c-e-f)+(b+c-e-f)^2.
\end{aligned}
\]
\par\noindent\hfill\textsc{Klein.\quad Schlesinger.}
\par\endgroup
